% Options for packages loaded elsewhere
\PassOptionsToPackage{unicode}{hyperref}
\PassOptionsToPackage{hyphens}{url}
\documentclass[
  ignorenonframetext,
]{beamer}
\newif\ifbibliography
\usepackage{pgfpages}
\setbeamertemplate{caption}[numbered]
\setbeamertemplate{caption label separator}{: }
\setbeamercolor{caption name}{fg=normal text.fg}
\beamertemplatenavigationsymbolsempty
% remove section numbering
\setbeamertemplate{part page}{
  \centering
  \begin{beamercolorbox}[sep=16pt,center]{part title}
    \usebeamerfont{part title}\insertpart\par
  \end{beamercolorbox}
}
\setbeamertemplate{section page}{
  \centering
  \begin{beamercolorbox}[sep=12pt,center]{section title}
    \usebeamerfont{section title}\insertsection\par
  \end{beamercolorbox}
}
\setbeamertemplate{subsection page}{
  \centering
  \begin{beamercolorbox}[sep=8pt,center]{subsection title}
    \usebeamerfont{subsection title}\insertsubsection\par
  \end{beamercolorbox}
}
% Prevent slide breaks in the middle of a paragraph
\widowpenalties 1 10000
\raggedbottom
\AtBeginPart{
  \frame{\partpage}
}
\AtBeginSection{
  \ifbibliography
  \else
    \frame{\sectionpage}
  \fi
}
\AtBeginSubsection{
  \frame{\subsectionpage}
}
\usepackage{iftex}
\ifPDFTeX
  \usepackage[T1]{fontenc}
  \usepackage[utf8]{inputenc}
  \usepackage{textcomp} % provide euro and other symbols
\else % if luatex or xetex
  \usepackage{unicode-math} % this also loads fontspec
  \defaultfontfeatures{Scale=MatchLowercase}
  \defaultfontfeatures[\rmfamily]{Ligatures=TeX,Scale=1}
\fi
\usepackage{lmodern}
\ifPDFTeX\else
  % xetex/luatex font selection
\fi
% Use upquote if available, for straight quotes in verbatim environments
\IfFileExists{upquote.sty}{\usepackage{upquote}}{}
\IfFileExists{microtype.sty}{% use microtype if available
  \usepackage[]{microtype}
  \UseMicrotypeSet[protrusion]{basicmath} % disable protrusion for tt fonts
}{}
\makeatletter
\@ifundefined{KOMAClassName}{% if non-KOMA class
  \IfFileExists{parskip.sty}{%
    \usepackage{parskip}
  }{% else
    \setlength{\parindent}{0pt}
    \setlength{\parskip}{6pt plus 2pt minus 1pt}}
}{% if KOMA class
  \KOMAoptions{parskip=half}}
\makeatother
\usepackage{color}
\usepackage{fancyvrb}
\newcommand{\VerbBar}{|}
\newcommand{\VERB}{\Verb[commandchars=\\\{\}]}
\DefineVerbatimEnvironment{Highlighting}{Verbatim}{commandchars=\\\{\}}
% Add ',fontsize=\small' for more characters per line
\newenvironment{Shaded}{}{}
\newcommand{\AlertTok}[1]{\textcolor[rgb]{1.00,0.00,0.00}{\textbf{#1}}}
\newcommand{\AnnotationTok}[1]{\textcolor[rgb]{0.38,0.63,0.69}{\textbf{\textit{#1}}}}
\newcommand{\AttributeTok}[1]{\textcolor[rgb]{0.49,0.56,0.16}{#1}}
\newcommand{\BaseNTok}[1]{\textcolor[rgb]{0.25,0.63,0.44}{#1}}
\newcommand{\BuiltInTok}[1]{\textcolor[rgb]{0.00,0.50,0.00}{#1}}
\newcommand{\CharTok}[1]{\textcolor[rgb]{0.25,0.44,0.63}{#1}}
\newcommand{\CommentTok}[1]{\textcolor[rgb]{0.38,0.63,0.69}{\textit{#1}}}
\newcommand{\CommentVarTok}[1]{\textcolor[rgb]{0.38,0.63,0.69}{\textbf{\textit{#1}}}}
\newcommand{\ConstantTok}[1]{\textcolor[rgb]{0.53,0.00,0.00}{#1}}
\newcommand{\ControlFlowTok}[1]{\textcolor[rgb]{0.00,0.44,0.13}{\textbf{#1}}}
\newcommand{\DataTypeTok}[1]{\textcolor[rgb]{0.56,0.13,0.00}{#1}}
\newcommand{\DecValTok}[1]{\textcolor[rgb]{0.25,0.63,0.44}{#1}}
\newcommand{\DocumentationTok}[1]{\textcolor[rgb]{0.73,0.13,0.13}{\textit{#1}}}
\newcommand{\ErrorTok}[1]{\textcolor[rgb]{1.00,0.00,0.00}{\textbf{#1}}}
\newcommand{\ExtensionTok}[1]{#1}
\newcommand{\FloatTok}[1]{\textcolor[rgb]{0.25,0.63,0.44}{#1}}
\newcommand{\FunctionTok}[1]{\textcolor[rgb]{0.02,0.16,0.49}{#1}}
\newcommand{\ImportTok}[1]{\textcolor[rgb]{0.00,0.50,0.00}{\textbf{#1}}}
\newcommand{\InformationTok}[1]{\textcolor[rgb]{0.38,0.63,0.69}{\textbf{\textit{#1}}}}
\newcommand{\KeywordTok}[1]{\textcolor[rgb]{0.00,0.44,0.13}{\textbf{#1}}}
\newcommand{\NormalTok}[1]{#1}
\newcommand{\OperatorTok}[1]{\textcolor[rgb]{0.40,0.40,0.40}{#1}}
\newcommand{\OtherTok}[1]{\textcolor[rgb]{0.00,0.44,0.13}{#1}}
\newcommand{\PreprocessorTok}[1]{\textcolor[rgb]{0.74,0.48,0.00}{#1}}
\newcommand{\RegionMarkerTok}[1]{#1}
\newcommand{\SpecialCharTok}[1]{\textcolor[rgb]{0.25,0.44,0.63}{#1}}
\newcommand{\SpecialStringTok}[1]{\textcolor[rgb]{0.73,0.40,0.53}{#1}}
\newcommand{\StringTok}[1]{\textcolor[rgb]{0.25,0.44,0.63}{#1}}
\newcommand{\VariableTok}[1]{\textcolor[rgb]{0.10,0.09,0.49}{#1}}
\newcommand{\VerbatimStringTok}[1]{\textcolor[rgb]{0.25,0.44,0.63}{#1}}
\newcommand{\WarningTok}[1]{\textcolor[rgb]{0.38,0.63,0.69}{\textbf{\textit{#1}}}}
\usepackage{longtable,booktabs,array}
\newcounter{none} % for unnumbered tables
\usepackage{calc} % for calculating minipage widths
\usepackage{caption}
% Make caption package work with longtable
\makeatletter
\def\fnum@table{\tablename~\thetable}
\makeatother
\setlength{\emergencystretch}{3em} % prevent overfull lines
\providecommand{\tightlist}{%
  \setlength{\itemsep}{0pt}\setlength{\parskip}{0pt}}
\usepackage{bookmark}
\IfFileExists{xurl.sty}{\usepackage{xurl}}{} % add URL line breaks if available
\urlstyle{same}
\hypersetup{
  hidelinks,
  pdfcreator={LaTeX via pandoc}}

\author{\texorpdfstring{}{}}
\date{}

\begin{document}

\subsection{Quality Assurance}\label{quality-assurance}

\begin{frame}[fragile]{Zero-Mock Testing Policy}
\protect\phantomsection\label{zero-mock-testing-policy}
All tests use real methods exclusively {[}@martin2008clean;
@meszaros2007xunit{]}. No \texttt{unittest.mock}, no \texttt{MagicMock},
no \texttt{patch} decorators. Tests that require external services
(Ollama, network) use \texttt{pytest.mark} markers for conditional
execution. The philosophical motivation---analogizing mock objects to
Simmons et al.'s \emph{researcher degrees of freedom}
{[}@simmons2011falsepositive{]} and the pre-registration remedy
{[}@nosek2018preregistration{]}---is developed fully in the
\hyperlink{the-zero-mock-tradeoff}{Zero-Mock Tradeoff} discussion. To
our knowledge, no prior research software engineering framework has
formalized a zero-mock policy as an \emph{architectural invariant
enforced by pipeline gates}, where mock usage is not merely discouraged
but structurally prevented from passing the build.

The following example, drawn from the infrastructure test suite,
illustrates zero-mock compliance:

\begin{Shaded}
\begin{Highlighting}[]
\KeywordTok{def}\NormalTok{ test\_discover\_infrastructure\_modules\_returns\_nonempty(tmp\_path):}
    \CommentTok{\# Real filesystem, real YAML parsing — no MagicMock anywhere}
\NormalTok{    modules }\OperatorTok{=}\NormalTok{ discover\_infrastructure\_modules(REPO\_ROOT)}
    \ControlFlowTok{assert} \BuiltInTok{len}\NormalTok{(modules) }\OperatorTok{\textgreater{}=} \DecValTok{8}           \CommentTok{\# actual subpackages on disk}
    \ControlFlowTok{assert} \BuiltInTok{any}\NormalTok{(m.name }\OperatorTok{==} \StringTok{"core"} \ControlFlowTok{for}\NormalTok{ m }\KeywordTok{in}\NormalTok{ modules)}
\end{Highlighting}
\end{Shaded}

This test exercises the real \texttt{discover\_infrastructure\_modules}
function against the real filesystem. There are no mock objects
substituting for the directory walk, no patched YAML parsers, and no
synthetic return values---the test passes only if the infrastructure
modules genuinely exist and are discoverable at their expected paths.
\end{frame}

\begin{frame}[fragile]{Coverage Thresholds}
\protect\phantomsection\label{coverage-thresholds}
The pipeline enforces two coverage tiers:

{\def\LTcaptype{none} % do not increment counter
\begin{longtable}[]{@{}
  >{\raggedright\arraybackslash}p{(\linewidth - 8\tabcolsep) * \real{0.1429}}
  >{\raggedright\arraybackslash}p{(\linewidth - 8\tabcolsep) * \real{0.1667}}
  >{\centering\arraybackslash}p{(\linewidth - 8\tabcolsep) * \real{0.2143}}
  >{\centering\arraybackslash}p{(\linewidth - 8\tabcolsep) * \real{0.2143}}
  >{\raggedright\arraybackslash}p{(\linewidth - 8\tabcolsep) * \real{0.2619}}@{}}
\toprule\noalign{}
\begin{minipage}[b]{\linewidth}\raggedright
Tier
\end{minipage} & \begin{minipage}[b]{\linewidth}\raggedright
Scope
\end{minipage} & \begin{minipage}[b]{\linewidth}\centering
Minimum
\end{minipage} & \begin{minipage}[b]{\linewidth}\centering
Current
\end{minipage} & \begin{minipage}[b]{\linewidth}\raggedright
Rationale
\end{minipage} \\
\midrule\noalign{}
\endhead
Project & \texttt{projects/*/src/} & 90\% & 90\%+ & Domain code must be
thoroughly validated \\
Infrastructure & \texttt{infrastructure/} & 60\% & 83\%+ & Broader
scope, some code unreachable in test \\
\bottomrule\noalign{}
\end{longtable}
}

These thresholds are enforced at Stage 01 of the pipeline. If project
test coverage falls below 90\%, the pipeline halts and refuses to
produce a PDF---ensuring that no published artifact is backed by
undertested source code.
\end{frame}

\begin{frame}[fragile]{Test Suite Composition}
\protect\phantomsection\label{test-suite-composition}
The repository maintains three test suites:

\begin{itemize}
\tightlist
\item
  \textbf{Infrastructure tests} (\texttt{tests/}): \textasciitilde3,083
  tests validating the 12 infrastructure subpackages, covering logging,
  rendering, validation, steganography, reporting, and LLM integration.
\item
  \textbf{Project tests} (\texttt{projects/*/tests/}): Per-project test
  suites validating domain-specific logic. Sizes vary from 39 tests
  (\texttt{code\_project}) to 505 tests
  (\texttt{act\_inf\_metaanalysis}).
\item
  \textbf{Integration tests}: Embedded within infrastructure tests,
  these exercise full pipeline stages against real manuscript inputs,
  validating end-to-end behavior from Markdown source to rendered PDF.
\end{itemize}
\end{frame}

\begin{frame}[fragile]{Visualization Standards}
\protect\phantomsection\label{visualization-standards}
All generated figures must meet accessibility requirements:

\begin{itemize}
\tightlist
\item
  Minimum 16pt font size for all text elements (the accessibility
  floor).
\item
  Colorblind-safe palettes (IBM Design / Wong palette) with
  high-contrast fallbacks.
\item
  150--300 DPI rendering for publication quality.
\item
  Descriptive axis labels and figure titles.
\item
  No reliance on color alone to convey information---redundant encoding
  via shape, pattern, or annotation is used where applicable.
\end{itemize}

These standards are validated by the \texttt{test\_architecture\_viz.py}
test suite, which verifies that generated figures exist, have non-zero
file size, and conform to expected output specifications. The 16pt font
floor ensures readability in both screen and print contexts, while the
DPI range balances file size against reproduction fidelity.
\end{frame}

\end{document}
